In the previous three sections, we have discussed extensively the theoretical methodologies to compute the energy of atomistic structures of materials, to sample the materials' energy landscapes, and to determine the physical properties of materials. Although these methodologies appear to have covered a wide range of applicative scenarios, for a long period of time before the adaptation of the big-data-driven, as well as the open-sourced, research software and data practices, computational materials science has primarily been the endeavour of individual research groups, each performing calculations on a specific (or a small set of) material(s) by applying the methodologies discussed above. There are pros and cons to this traditional approach. Fundamentally, this was limited by how many of the computational resources were accessible to the individual research group. In the lack of community-based best practices for performing materials simulations, each researcher was required to focus on optimising the simulation parameters for a single structure in their own interests, that is often benchmarked against some known experimental data. While this promotes each researcher to exploit the scientific depth to the material system he/she is interested in, it also leads to significant limitations in the data thus generated. 

For instances, firstly, the generalisability and validity of the simulation methodologies adopted in each isolated study become questionable, meaning it is difficult to judge whether the same method will work equally well to other materials system alike, or whether the agreement with (mostly experimental) observation is incidental, that could have led to biased or even incorrect conclusions in the underlying mechanisms behind the physical or chemical properties investigated in each study. As mentioned before, in the lack of systematic benchmarks, selecting computational methods based on `popularity' not appropriateness or validity, such as PBE for solid-state and B3LYP\cite{kruse2012standard} for molecular systems thus had been, and to some extent in the present days, remains to be the `default' approach of density functional selection(s) in DFT-based materials studies. Secondly, research groups tend to develop their own set of computer scripts/programs to derive further physical properties by post-processing and analysing the DFT results, often leads to low reproducibility of similar  works across the research community. This heterogeneity in research methodologies together with the lack of open-sourcing practice in computational materials research naturally lead to challenges such as to perform statistical analysis on computational results to identify systematic structure-property relationships, as well as outliers amongst a group of materials, which significantly limited our capabilities to generate, not just qualitatively correct but also quantitatively useful insights, to guide the materials discoveries. 

On the experimental side, traditional trial-and-error approaches often lead to slow and labour-intensive cycles in identifying materials with desired properties. Over the past decades, advances in computational power, especially the emergence of GPU-accelerated materials simulations softwares, have fundamentally transformed this process. Large-scale computational pre-screenings based on either experimetally or computationally generated data, that are increasingly enhanced by ML/AI algorithms, can now efficiently narrow down candidate materials and substantially raise the success rate of experimental synthesis and discovery. These capabilities have become a major driving force behind the widespread adoption of high-throughput computational materials research across the community. 

For examples, \citet{raccuglia2016machine} applied ML algorithms to learn from a database containing 3,955 reactions, including failed ones, to guide the synthesis of new compounds that showed a significant increase of success rate to 89 \%. To facilitate the designs of molecular structures, \citet{ahneman2018predicting} applied ML algorithms and the Pd-catalysed Buchwald-Hartwig cross-coupling of aryl halides with 4-methylaniline under different inhibitory additives as examples, showcased how such theory-guided experimental approach was successful in minimising the side reactions, and enhance the stability of the catalyst towards the successful synthesis of the desired product.  Measuring the liftime of lithium-ion batteries is important for their commercialisation but also a time-consuming process. To accelerate the process of life-cycle measurements, \citet{severson2019data} devised ML approaches to extract features in the discharge voltage curves from early cycles, from which they were able to predict and classify the cells by life cycles. Similarly, \citet{zahrt2019prediction} applied ML-guided molecular design of organocatalysts to show how the reaction yields and stereoselectivities could be significantly enhanced. Despite the limitless choices in the structures of organic molecules, only less than 2000 organic photovoltaic donor molecules have been synthesised and tested in the photovoltaic cells, since organo-photovoltaics were first demonstrated in 1973,\cite{ghosh1973rectification} this reflects the long research cycles that involves molecular design, delicate control of the chemical synthesis conditions, laborious purifications and crystallisation processes as well as the challenges behind device fabrications. To overcome these challenges, \citet{sun2019machine} applied ML approach to extract the structure-property relationships from over 1700 experimentally reported chemicals, which led to the successful discovery of 10 newly designed donor molecules. 

The successes behind these works, which nowadays are commonly referred to as the \emph{data-driven} approach in materials discoveries or \emph{materials informatics}, are largely backed by the unique advantages for it to algorithmically extract hidden relationships among the chemical compositions, structures and the properties behind materials over a large materials dataset in a non-biased way, overcoming the limitations behind human intuitions. This has significantly enhanced the efficiency for us the navigate the materials spaces to find better candidates, identify potential pitfalls in the experiments, thus allowing us to better prioritise the compositions and conditions to synthesise new materials, ultimately shorten the materials development lifecycles. 

The paradigm shift from individuals performing isolated computational studies towards the global and communal effort in assembling large-scale material databases via high-throughput computations,\cite{jain2013commentary,saal2013materials,choudhary2020high,draxl2019nomad,haas2009calculation,gjerding2021recent,chanussot2021open,tran2023open} was not only to harness the continuously increasing available computational powers, but most importantly, to address and overcome the limitations mentioned above on the shortfalls behind isolated research efforts, as well as to further capitalise the success of data-driven material discoveries. Such a paradigm shift, together with the widely applied methods of structure predictions/energy landscape sampling as discussed in \cref{chapter:sampling_methods} have significantly transformed the field of materials science by allowing us to establish a holistic view of materials across both of its composition and structural spaces [\cref{fig:energy_landscapes_hierachical}].  In this section, we shall provide a brief account on  the development of materials databases that is enabled by high-throughput computations backed by efficient automated workflow management.
